\documentclass[12pt]{article}
\usepackage[T1]{fontenc}
\usepackage[utf8]{inputenc}
\usepackage{lmodern}
\usepackage{textcomp}
\usepackage{enumitem}
\usepackage{geometry}
\geometry{margin=1in}
\begin{document}
\begin{center}
Answer Key - Psychology - Graduate Level Assessment 12 \
\small (Answers are ordered exactly as in the question paper.)
\end{center}

\section*{Multiple Choice}
\begin{enumerate}[label=\arabic*.]
\item \textbf{Answer:} A
\\ \textit{Options:} A) Attractiveness can be precisely measured and quantified in all individuals.; B) We fully understand the factors that determine human attractiveness.; C) Attractiveness is a trait that humans are born understanding completely.; D) Human attractiveness is determined by a single genetic marker.; E) The principles of attractiveness in humans have been fully mapped out by neuroscience.
\\ \textit{Note:} The correct answer highlights that the quantifiable nature of attractiveness in fish allows for more precise and consistent measurements across individuals, whereas human attraction is influenced by complex, subjective factors that are more difficult to quantify.
\item \textbf{Answer:} A
\\ \textit{Options:} A) During adolescence only; B) From birth to late childhood only; C) During infancy only; D) From birth to early adulthood only; E) From infancy through adolescence only
\\ \textit{Note:} Piaget's theory describes the final stage of cognitive development, the formal operational stage, which occurs during adolescence and involves the development of abstract reasoning and hypothetical thinking.
\item \textbf{Answer:} A
\\ \textit{Options:} A) Bandura's social learning theory.; B) Skinner's operant conditioning theory.; C) Maslow's hierarchy of needs theory.; D) Beck's cognitive-behavioral therapy.; E) Pavlov's classical conditioning theory.
\\ \textit{Note:} The correct answer is Bandura's social learning theory because it emphasizes the role of social context and observational learning in shaping behavior, suggesting that a child's misbehavior can serve specific social goals such as seeking attention or asserting power.
\item \textbf{Answer:} A
\\ \textit{Options:} A) Childhood experiences; B) Astrological signs and birth months; C) Educational background and intelligence quotient (IQ); D) Exercise frequency and body mass index (BMI); E) Use of technology and media consumption
\\ \textit{Note:} Psychologists generally accept that childhood experiences, including trauma and adverse events, significantly influence the development of functional psychoses by shaping an individual's psychological resilience and coping mechanisms.
\item \textbf{Answer:} E
\\ \textit{Options:} A) Negligence in supervision; B) Incompetence; C) Misuse of authority; D) Sexual harassment; E) Lack of timely feedback
\\ \textit{Note:} The correct answer is "Lack of timely feedback" because failure to provide prompt and constructive evaluation can lead to misunderstandings and ethical dilemmas regarding the supervision process, ultimately resulting in complaints against supervisors.
\end{enumerate}

\section*{True / False}
\begin{enumerate}[label=\arabic*.]
\setcounter{enumi}{5}
\item \textbf{Answer:} True
\\ \textit{Options:} A) True; B) False
\\ \textit{Note:} Surprise is not classified as a primary facial expression in psychological research because the primary expressions typically identified are happiness, sadness, anger, fear, disgust, and contempt, while surprise is often considered a secondary or mixed expression that can occur in conjunction with others.
\item \textbf{Answer:} False
\\ \textit{Options:} A) True; B) False
\\ \textit{Note:} Time-out procedures are based on operant conditioning, where the removal of a positive stimulus follows a specific behavior to reduce its occurrence, rather than classical conditioning, which involves associating an involuntary response with a stimulus.
\item \textbf{Answer:} True
\\ \textit{Options:} A) True; B) False
\\ \textit{Note:} Stimulus satiation is the psychological phenomenon where the repeated exposure to a stimulus leads to a decrease in its appeal or desirability, thereby confirming that the statement is true.
\item \textbf{Answer:} False
\\ \textit{Options:} A) True; B) False
\\ \textit{Note:} Strong organizational culture is correlated with high levels of job commitment in the workplace because a positive culture fosters employee engagement, loyalty, and alignment with organizational values.
\item \textbf{Answer:} False
\\ \textit{Options:} A) True; B) False
\\ \textit{Note:} The Whorfian hypothesis, also known as linguistic relativity, posits that language influences thought and perception rather than asserting that all languages share identical phonetic structures or universal tone systems.
\end{enumerate}

\section*{Long Form Response}
\begin{enumerate}[label=\arabic*.]
\setcounter{enumi}{10}
\item \textbf{Answer:} Type II error occurs when a researcher fails to reject a null hypothesis that is false, leading to a missed opportunity to identify a true effect. Statistical power, defined as the probability of correctly rejecting a false null hypothesis, is influenced by several factors, with sample size being a critical component. Specifically, as sample size increases, the statistical power of a test generally increases, enhancing the likelihood of detecting true effects. Conversely, if the sample size is small, the power decreases, making it more difficult to identify significant results, which can lead to Type II errors. This relationship underscores the importance of adequately powered studies in psychology to ensure that genuine effects are not overlooked.
\\ \textit{Note:} correct because it accurately defines Type II error as failing to reject a false null hypothesis and explains that increasing sample size enhances statistical power, thereby improving the test's ability to detect true effects.
\item \textbf{Answer:} The basal ganglia play a crucial role in the Cannon-Bard theory of emotion by mediating the integration of emotional stimuli with motor responses. This theory posits that emotional experiences occur simultaneously with physiological reactions, challenging the earlier James-Lange theory, which suggested that emotions follow bodily responses. The basal ganglia, particularly through their connections with the limbic system, facilitate the expression of emotions by influencing movement and motivation, thereby reinforcing emotional responses. Understanding the function of the basal ganglia in emotional processing highlights the complexity of how emotions are not only felt but also expressed behaviorally, providing insights into disorders where these processes may be disrupted. Consequently, this reinforces the notion that emotions arise from a network of brain regions working in concert rather than in a linear sequence.
\\ \textit{Note:} correct because it accurately identifies the basal ganglia's role in the Cannon-Bard theory as a mediator of emotional and motor responses, emphasizing the simultaneous experience of emotions and physiological reactions, which is central to understanding emotional processing in the brain.
\end{enumerate}

\vfill
\noindent\textit{End of Answer Key}
\end{document}